\documentclass[12pt]{article}
\usepackage{amsmath,amssymb,amsfonts}
\usepackage[margin=1in]{geometry}

\begin{document}

\section*{Problem 5.16}

\textbf{Problem:} Prove that the complete set of orthonormal functions on the unit sphere, $Y_{\ell m}$, is unique up to a phase factor. Use induction on $\ell$ for $m = \ell$ or $m = -\ell$.

\bigskip
\textbf{Solution:}

The spherical harmonics $Y_{\ell m}(\theta, \phi)$ form a complete orthonormal set on the unit sphere. We want to show that if $\{Y_{\ell m}\}$ and $\{Y'_{\ell m}\}$ are two such complete orthonormal sets satisfying the same eigenvalue equations, then they differ at most by a phase factor:
\[
Y'_{\ell m}(\theta, \phi) = e^{i\alpha_{\ell m}} Y_{\ell m}(\theta, \phi).
\]

\subsection*{Step 1: Uniqueness of the $\phi$-dependence}

The spherical harmonics satisfy:
\[
L_z Y_{\ell m} = -i\frac{\partial}{\partial\phi}Y_{\ell m} = m Y_{\ell m}.
\]

This first-order ODE in $\phi$ has a unique solution (up to a function of $\theta$):
\[
Y_{\ell m}(\theta, \phi) = \Theta_{\ell m}(\theta)\,e^{im\phi}.
\]

Any other solution must have the same $e^{im\phi}$ dependence, so:
\[
Y'_{\ell m}(\theta, \phi) = \Theta'_{\ell m}(\theta)\,e^{im\phi}.
\]

\subsection*{Step 2: Uniqueness of $\Theta_{\ell m}(\theta)$ up to phase}

The $\theta$-dependent part satisfies the associated Legendre equation:
\[
\frac{1}{\sin\theta}\frac{d}{d\theta}\left(\sin\theta\frac{d\Theta}{d\theta}\right) + \left[\ell(\ell+1) - \frac{m^2}{\sin^2\theta}\right]\Theta = 0.
\]

This second-order ODE has two linearly independent solutions, but only one is regular at both $\theta = 0$ and $\theta = \pi$: the associated Legendre function $P_\ell^m(\cos\theta)$. Therefore:
\[
\Theta_{\ell m}(\theta) = c_{\ell m}\,P_\ell^m(\cos\theta), \qquad \Theta'_{\ell m}(\theta) = c'_{\ell m}\,P_\ell^m(\cos\theta),
\]
where $c_{\ell m}$ and $c'_{\ell m}$ are constants.

\subsection*{Step 3: Normalization fixes the magnitude}

The normalization condition $\int |Y_{\ell m}|^2\,d\Omega = 1$ fixes $|c_{\ell m}|$. Similarly for $|c'_{\ell m}|$. Thus $|c_{\ell m}| = |c'_{\ell m}|$, and:
\[
c'_{\ell m} = e^{i\alpha_{\ell m}}\,c_{\ell m}
\]
for some real phase $\alpha_{\ell m}$.

\subsection*{Step 4: Induction argument for $m = \pm\ell$}

\textbf{Base case:} For $\ell = 0$, $Y_{00} = 1/\sqrt{4\pi}$ is unique (real and positive by convention), so $Y'_{00} = e^{i\alpha_{00}}Y_{00}$.

\textbf{Inductive step:} Assume uniqueness up to phase for spherical harmonics up to $\ell - 1$. For $\ell$, consider $m = \ell$. The associated Legendre function $P_\ell^\ell(\cos\theta) \propto \sin^\ell\theta$, so:
\[
Y_{\ell\ell}(\theta,\phi) = c_\ell\,\sin^\ell\theta\,e^{i\ell\phi}.
\]

Normalization fixes $|c_\ell|$. The remaining freedom is a single phase factor. The ladder operators $L_\pm = L_x \pm iL_y$ then generate all other $m$ values:
\[
Y_{\ell m} \propto L_-^{\ell-m}\,Y_{\ell\ell}.
\]

Since $L_-$ is a fixed differential operator, $Y_{\ell m}$ is determined up to the same phase factor as $Y_{\ell\ell}$.

\subsection*{Conclusion}

\[
\boxed{Y'_{\ell m} = e^{i\alpha_{\ell m}}\,Y_{\ell m},}
\]
and the complete set of orthonormal functions on the unit sphere is unique up to phase factors. With the Condon--Shortley phase convention, the phase is completely fixed.

\end{document}
